\documentclass[12pt]{article}
\usepackage{resource/shortex}
\input{resource/teaching}

%--------------------
\begin{document}

\begin{center}
\textbf{\Large Project II Rubric} 
\end{center}

\grades

\subsection*{Rubric}


\questiongrades

\noindent
\section*{General Rubric}
\subsubsection*{Data Description and Introduction}
\normalsize
\begin{tabular}{lp{.75\textwidth}}
\toprule
\textbf{Mark} & \textbf{Description} \\
\hline
\triAssessments & \textbf{(1)} Reasoning and motivation are clearly stated \\
\triAssessments & \textbf{(2)}  Dataset is clearly introduced  \\
% \triAssessments & \textbf{(3)} Variables of interest (numerical and categorical) are well defined \\
\triAssessments & \textbf{(3)} Sentences are well structured, easy to follow, and includes smooth transitions between sections \\
\triAssessments & \textbf{(4)} Visualizations and tables are integrated clearly and contribute meaningfully to the analysis \\
\triAssessments & \textbf{(5)} R code is clean, well commented, reproducible, and submitted in the correct format \\
\triAssessments & \textbf{(6)} Proper citations and references are included for datasets and any external resources \\
\bottomrule
\end{tabular}
\vspace{1em} 


\subsubsection*{Categorical Data Analysis}
\normalsize
\begin{tabular}{lp{.75\textwidth}}
\toprule
\textbf{Mark} & \textbf{Description} \\
\hline
\triAssessments & \textbf{(1)} Provides appropriate exploratory summaries of contingency tables of categorical variables \\
\triAssessments & \textbf{(2)} Includes at lesast one visualizations (bar plot or pie chart) with proper labels and titles \\
\triAssessments & \textbf{(3)} States and checks conditions necessary for inference \\
\triAssessments & \textbf{(4)} Formulates correct hypotheses for the test \\
\triAssessments & \textbf{(5)} Calculates the correct test statistic and p-value accurately \\
\triAssessments & \textbf{(6)} Provides logical interpretation and conclusion in context \\
\bottomrule
\end{tabular}
\vspace{1em}


\subsubsection*{Numerical Data Analysis}
\normalsize
\begin{tabular}{lp{.75\textwidth}}
\toprule
\textbf{Mark} & \textbf{Description} \\
\hline
\triAssessments & \textbf{(1)} Summarizes numerical data properly (mean, median, sd, IQR, etc.) \\
\triAssessments & \textbf{(2)} Includes at least one visualizations (histograms, boxplots, scatterplots) with clear labels \\
\triAssessments & \textbf{(3)} Checks and states conditions for inference methods used  \\
\triAssessments & \textbf{(4)} Correctly states hypotheses for numerical inference \\
\triAssessments & \textbf{(5)} Computes appropriate test statistic and p-value accurately \\
\triAssessments & \textbf{(6)} Provides clear conclusion with interpretation in data context \\
\bottomrule
\end{tabular}
\vspace{1em}



\section*{Individual Contribution Rubric}
\normalsize
\begin{tabular}{lp{.75\textwidth}}
\toprule
\textbf{Mark} & \textbf{Description} \\
\hline
\triAssessments & \textbf{(1)} Completed assigned individual tasks thoroughly \\
\triAssessments & \textbf{(2)} Actively participated in group collaboration and communication \\
\bottomrule
\end{tabular}
\vspace{1em}

\subsection*{Grading Standards}
\begin{itemize}
\item \textbf{E: Excellent}: Satisfactory with exceptional creativity or insight
\item \textbf{S: Satisfactory}: Any combination of $\checkmark$ and $\checkmark\!-$ 
\item \textbf{R: Revisions Needed}: Any \ding{55} on any criterion
\item \textbf{ N: Not Assessable}: Cannot grade due to too many missing components
\end{itemize}


\end{document}